% 第19章 局限性摘要
\section{局限性摘要与反方证据}\label{sec:limitations}

一份诚实的综述，应该把自己的边界写到别人能攻击的程度。这一章的写法是：先列出本文自身的局限，再把文献里与本文判断相反的证据摆出来，最后给出修正后的判断。

\subsection{本综述自身的局限}

\begin{enumerate}
  \item \textbf{检索覆盖偏差}：检索以英文文献与 arXiv 为主，可能低估中文与非英语圈的实证工作；对期刊的覆盖也偏向金融与机器学习主流刊物。
  \item \textbf{私有数据不可见}：对冲基金与做市商的真实业绩不可得，本文关于“学术结论与实盘落差”的判断，依据的是公开证据的衰减模式，而非直采的实盘数据。
  \item \textbf{参数细节需全文复核}：证据表中若只依据摘要与元数据，样本区间、阈值、成本模型等细节仍可能被简化，已逐条标注核验状态。
  \item \textbf{方法演进的时效}：大模型与生成式模型进展极快，本文对 LLM 类工作的保守评价，可能随前向验证证据的积累而需要修正。
  \item \textbf{篇幅取舍}：全景的双轴框架降低了单格的深度；对某一具体方法族想深挖的读者，仍需回到原始论文。
\end{enumerate}

\subsection{反方证据}

\textbf{反方一：AI 的增量可能被高估。}这一派（以有效市场传统的实证为据）认为，考虑交易成本、流动性与容量后，绝大多数所谓 AI 策略的超额收益会消失；把机器学习用在截面上，本质上只是“更高效地挖已知因子”，而非创造新信息。支持这条的间接证据是：因子被发表后收益衰减\cite{mcleanpontiff2016}、大批因子无法复现\cite{jensen2023}。本文承认这条批评的分量，这也是为什么全篇把“扣成本后的样本外”作为唯一有效判据。

\textbf{反方二：AI 的增量可能被低估。}“复杂性红利”\cite{kellymalamudzhou2024} 与部分端到端\cite{costa2022}、结构注入\cite{chenpelgerzhu2024} 的结果显示，在正确定义的问题上，更复杂的模型确有增量。这一派的潜台词是：传统实证低估了非线性与高维交互的价值，因为它的工具够不到那里。本文对此保持开放，但坚持一点——\textbf{这类增量必须在统一口径、扣成本、经多重检验后仍然存在，才算数}。

\textbf{反方三：LLM 智能体是被低估的范式。}乐观者认为，LLM 带来了“理解世界”的能力，而不仅是“拟合数据”，因此智能体交易的潜力被本文的保守评价低估。本文的回应是：\textbf{潜力不等于证据。}在“模型知道答案”这一污染被系统性排除之前，任何乐观判断都缺少认识论基础；本文并不否认范式潜力，只是拒绝在没有前向证据前给它定高的证据等级。

\subsection{本文的立场，修正后}

综合正反两方，本文修正后的判断是三条：

\begin{enumerate}
  \item \textbf{AI 提升的是“转化效率”，其可交易增量应随市场效率递减。}需要说明的是，“AI 不改变利润总量”这一说法本身不可观测、也不可证伪，本文不把它当作结论，而给出一个\emph{可检验}的弱化形式：\textbf{在统一口径（同数据、同成本、同评测协议）下，AI 类方法相对“最强简单基线”的成本后净增量，随市场效率的提升而递减。}检验设计是：取三个成熟度不同的市场（如美股大盘、A 股主板、加密永续），在同一管线上估计相对净增量，再看其排序是否与常见的效率代理（买卖价差、机构持股占比、信息半衰期）呈同序。排序成立则该命题获得支持；不成立则本文关于“效率上限”的论证需要修正。这正是第 \ref{sec:future} 章路线图 Phase 3 的任务之一。
  \item \textbf{越靠近摩擦的环节，AI 的增量越可靠；越靠近纯预测的环节，增量越不确定。}因为摩擦是可以被工程降低的客观量，而预测增量必须与市场效率上限竞争。这与“执行/做市类工作复现性好、选股类工作复现性差”的经验一致。
  \item \textbf{当前文献最大的问题不是模型，而是评测。}因此对读者而言，最高性价比的动作不是追逐更新奇的模型，而是先把口径与纪律立起来——这也是本文把 P0--P2 放在 P3 之前的原因。
\end{enumerate}

本章要记住的一句话是：\textbf{一份有用的综述不该让读者更兴奋，而该让读者更清醒——如果读完它你既没有高估 AI 也没有轻视它，那它就达到了目的。}
